% !TEX root = ../main.tex
\section{Markedssvikt og regulering}\label{thr:market-failure-regulation}

\paragraph{Definisjon.} \emph{Offentlig indgriben} (regulering) er hensigtsmæssig når markedet selv ikke oppnår en effisient allokering -- dvs.\ ved \emph{markedssvikt}. De fire klassiske argumentene er: (1) naturlig monopol, (2) eksternaliteter, (3) offentlige goder, og (4) informasjonsasymmetri. Regulering søker å korrigere priser, kvantiteter eller atferd slik at utfallet nærmer seg det en perfekt-konkurranse-benchmark ville gitt. [SLIDE]

\subsection*{Forklaring}

BSZ kap.\ 21 presenterer fire forhold der markedet feiler og regulering kan forbedre utfallet: [SLIDE/BOK]

\textbf{1. Naturlig monopol.} Produksjon med fallende gjennomsnittskostnad (stordrift, høye faste kostnader, lave variable) gjør at én produsent er billigst -- men uten konkurranse setter monopolisten $P>MC$ og produserer for lite. Regulering: \emph{prisregulering} ($P=AC$ eller rate-of-return-regulering), konsesjonskrav, offentlig eierskap. Eksempler: strømnett, vannforsyning, jernbaneskinner. [SLIDE]

\textbf{2. Eksternaliteter.} Aktiviteter med positive eller negative konsekvenser for tredjepart som ikke prises av markedet. Se \thref{externalities}{eksternaliteter}. Negative eks.: forurensning (overproduksjon); positive eks.: utdanning, FoU (underproduksjon). Regulering: Pigou-skatt/subsidie, utslippskvoter, standarder. [SLIDE]

\textbf{3. Offentlige goder.} Goder som er \emph{ikke-rivaliserende} (min bruk reduserer ikke din) og \emph{ikke-ekskluderbare} (vanskelig å hindre gratispassasjerer). Free-riding gjør at markedet underproduserer. Regulering: offentlig tilbud, finansiering via skatt. Eksempler: forsvar, grunnforskning, fyrtårn. [BOK]

\textbf{4. Informasjonsasymmetri.} Selger vet mer om kvalitet enn kjøper (\thref{adverse-selection}{adverse selection}), eller aktør skjuler handling (\thref{moral-hazard}{moral hazard}). Markedet kan kollapse (Akerlof lemons). Regulering: obligatorisk informasjonsplikt (regnskap, prospekt), minstestandarder, sertifisering, forsikringskrav. [BOK]

\textbf{Forutsætninger for at regulering er hensigtsmæssig:}
\begin{itemize}\setlength{\itemsep}{1pt}
\item Markedssvikten er \emph{signifikant} -- reguleringsomkostningene må være lavere enn dødvektstapet.
\item Regulatoren har \emph{tilstrekkelig informasjon} til å sette korrekte priser/kvantiteter.
\item Reguleringsregimet er \emph{ikke kapert} av bransjen selv (regulatory capture).
\end{itemize}

\subsection*{Kjerneforutsetninger}
\begin{itemize}
\item Perfekt konkurranse er benchmark for effisient allokering.
\item Markedssvikt er kilden til problemet, ikke distribusjonspolitikk alene.
\item Regulatoren er (ideelt) benevolent og velinformert.
\item Regulering har selv kostnader (byråkrati, feilinformasjon, lobbying).
\end{itemize}

\subsection*{Muligheter}
\begin{itemize}
\item Gir normativt rammeverk for når staten bør gripe inn.
\item Forklarer regulering av energi, telekom, finans, helse, miljø.
\item Kobler til konkrete reguleringsverktøy: priskontroll, utslippskvoter, tilgangsregulering.
\end{itemize}

\subsection*{Begrensninger og kritikk}
\begin{itemize}
\item \textbf{Government failure:} Regulatoren kan mangle informasjon, være byråkratisk, eller bli korrumpert.
\item \textbf{Regulatory capture} (Stigler 1971): bransjen påvirker regulatoren til å sette regler som beskytter incumbents.
\item \textbf{Reguleringsarbitrasje:} Aktører tilpasser seg reguleringen uten å endre reell atferd.
\item \textbf{Dynamiske effekter:} Regulering kan hemme innovasjon og entry -- f.eks.\ rate-of-return-regulering gir Averch-Johnson-overinvestering i kapital.
\item \textbf{Public choice:} Politikere optimerer stemmer, ikke velferd -- regulering reflekterer politisk støtte-maksimering, ikke markedssvikt-korreksjon. Se \thref{regulated-pricing-model}{regulated pricing-modellen}.
\end{itemize}

\subsection*{Modell / illustrasjon}

\begin{figure}[H]
  \centering
  \includegraphics[width=0.35\linewidth]{BSZ_kap_6_2026/by_slide/slide_003/BSZ_kap_6_2026_slide_003_asset_01_image2.png}
  \caption{Markedsstrukturer fra perfekt konkurranse til monopol. Regulering griper inn der markedsmakt gir $P > MC$. (BSZ kap.\ 6, slide 3) [SLIDE]}
\end{figure}

\begin{center}
\renewcommand{\arraystretch}{1.15}
\begin{tabular}{@{}p{3.5cm}p{3.8cm}p{3.5cm}p{3.5cm}@{}}
\toprule
\textbf{Markedssvikt} & \textbf{Problem} & \textbf{Reguleringsverktøy} & \textbf{Eksempel} \\
\midrule
Naturlig monopol & $P>MC$, underproduksjon & Priskontroll, konsesjoner & Strømnett, vannverk \\
Eksternaliteter & Over-/underproduksjon & Pigou-skatt, kvoter & CO$_2$-avgift, ETS \\
Offentlige goder & Free-riding, under\-produksjon & Offentlig tilbud, skatt & Forsvar, FoU \\
Informasjonsasymmetri & Lemons, adverse selection & Informasjonsplikt, standarder & Regnskapskrav, FDA \\
\bottomrule
\end{tabular}
\end{center}
\textit{Tabell: De fire markedssviktene og korresponderende reguleringsverktøy. [SLIDE/BOK]}

\subsection*{Eksempler}
\begin{itemize}\setlength{\itemsep}{2pt}
\item \textbf{Norsk strømnett (Statnett):} Naturlig monopol regulert av NVE med inntektsramme-modell (rate-of-return-regulering). Sikrer lav pris men risikerer overinvestering (Averch-Johnson). [EKS]
\item \textbf{EU Emissions Trading System (ETS):} Regulering av CO$_2$-eksternalitet via kvotehandel -- Pigou-mekanisme. Prisen ($\sim$50--80 EUR/tonn) internaliserer deler av den negative eksternaliteten. [EKS]
\item \textbf{Mærsk Power-to-X og subsidier.} PtX mottar offentlig støtte fordi grønt hydrogen produserer positive eksternaliteter (klimareduksjon) som markedet underpriser. Subsidiering er reguleringsresponsen. [CASE]
\end{itemize}

\subsection*{Kobling til søyle og andre teorier}
Søyle: \SOne\ (regulatoren allokerer beslutningsrettigheter mellom stat og marked). Relaterte: \thref{externalities}{Eksternaliteter}, \thref{adverse-selection}{Adverse selection}, \thref{moral-hazard}{Moral hazard}, \thref{regulated-pricing-model}{Regulated pricing-modellen}, \thref{vertical-integration-market-power}{VI og market power}.

\paragraph{Brukt i spørsmål:} Q13.
